\subsubsection{Subproceso 2: Selección y Registro de Cursos}



El Subproceso 2 arranca cuando el estudiante —que ya tiene su turno asignado en el primer paso— entra a la Intranet de la DIRCE dentro del horario que le tocó. Al iniciar sesión con su correo institucional (@uni.pe), el sistema valida que todo esté en orden y le muestra únicamente los cursos de su plan de estudios que le corresponden: aquellos cuyos requisitos ya aprobó y que están abiertos para este ciclo.

Luego, el alumno revisa las opciones y arma su horario eligiendo los cursos, secciones y grupos que prefiere, teniendo en cuenta la regla de priorizar los cursos pendientes o de riesgo. Mientras va armando su horario, la plataforma revisa al instante dos cosas: que no rebase el límite de créditos permitido según su promedio y que no se crucen las horas de clase. Si algo no cumple con las reglas, el sistema lo bloquea y le avisa para que haga los cambios necesarios.

Una vez que el estudiante está conforme con su horario y cumple con todas las normas, hace clic en registrar y enviar su propuesta. En ese momento, el sistema SIGA-ORCE aparta de forma automática los cupos en cada sección elegida y genera la Pre-Ficha de Matrícula Registrada, un comprobante provisional que resume los cursos, horarios y créditos del semestre. Con este documento se cierra el Subproceso 2 y la pre-ficha pasa al Subproceso 3.A para su revisión.

